
\section{Background}

Before discussing our method, we briefly introduce the background for IL, including learning from demonstrations (LfD), learning from observations (LfO), online IL, offline IL, and cross-domain settings in Section~\ref{sec:imitation-learning}, and introduce the hindsight information matching in Section~\ref{sec:hindsight-infomation-matching}. 

\subsection{Imitation Learning}
\label{sec:imitation-learning}

Consider a control task formulated as a discrete-time Markov decision process (MDP)\footnote{In this paper, we use environment and MDP interchangeably, and use state and observation interchangeably.} 
$\mathcal{M} = \{ \mathcal{S}, \mathcal{A}, \mathcal{T}, r, \gamma, p_0 \}$, where $\mathcal{S}$ is the state (observation) space, $\mathcal{A}$ is the action space, $\mathcal{T}: \mathcal{S} \times \mathcal{A} \times \mathcal{S} \to \mathbb{R}$ is the transition dynamics function, $r: \mathcal{S} \times \mathcal{A} \to \mathbb{R}$ is the reward function, $\gamma$ is the discount factor, and $p_0$ is the distribution of initial states. The goal in a reinforcement learning (RL) control task is to learn a policy $\pi_\theta(\ba|\bs)$ maximizing the expected sum of discounted rewards  $\mathbb{E}_{\pi_\theta(\btau)}\left[ \sum_{t=0}^{T-1} \gamma^t r(\bs_t, \ba_t) \right]$, where $\btau := \{\bs_0, \ba_0, \cdots, \bs_{T-1}, \ba_{T-1}\}$ denotes the trajectory and the generated trajectory distribution $\pi_\theta(\btau) = p_0(\bs_0)\pi_\theta(\ba_0|\bs_0)\prod_{t=1}^{T-1}\pi_\theta(\ba_t|\bs_t)\mathcal{T}(\bs_t|\bs_{t-1}, \ba_{t-1})$.
%\begin{align}\label{eq:generate-tau}
%	p_0(\bs_0)\pi_\theta(\ba_0|\bs_0)\prod_{t=1}^{T-1}\pi_\theta(\ba_t|\bs_t)T(\bs_t|\bs_{t-1}, \ba_{t-1}). 
%\end{align}

In IL, the ground truth reward function (\ie $r$ in $\mathcal{M}$) is not observed. Instead, we have access to a set of demonstrations (or observations) $\{\btau|\btau\sim\pi_E(\btau)\}$ that are collected by an unknown expert policy $\pi_E(\ba|\bs)$. 
%(Here we use the tilde $\tbtau$ to emphasize that trajectory $\tbtau$ may contain only observations/states, \ie actions in $\tbtau$ are absent.)
The goal of IL tasks is to recover a policy that matches the corresponding expert policy. From the mathematical perspective, IL achieves the plain expert matching objective by minimizing the divergence of trajectory distributions between the learner and the expert:  
\begin{align}\label{eq:min-il}
	\min_{\pi_\theta} \  D(\pi_\theta(\btau), \pi_E(\btau)), 
\end{align}
where $D$ is a distance measure. 
Meanwhile, we emphasize that the given expert data $\{\btau|\btau\sim\pi_E(\btau)\}$ may not contain the corresponding expert actions. Thus, in this work, we consider two IL cases where the given expert data $\btau$ consists of a set of state-action demonstrations $\{(\bs_t, \ba_t, \bs_{t+1})\}$ ({learning from demonstrations, {{LfD}}}), as well as a set of state-only transitions $\{(\bs_t, \bs_{t+1})\}$ ({learning from observations, {{LfO}}}). 
\textit{When it is clear from context, we abuse notation $\pi_E(\btau)$ to denote both demonstrations in LfD and observations in LfO for simplicity.} 

% To optimize Equation~\ref{eq:min-il}, typical LfD methods propose to compute the distribution discrepancy over the state-action occupancy measure $\rho_{\pi_\theta}(\bs, \ba)$, leading to
% \begin{align}%\label{eq:min-il-lfd}
% 	D(\pi_\theta(\btau), \pi_E(\btau)) \to D(\rho_{\pi_\theta}(\bs, \ba), \rho_{\pi_E}(\bs, \ba)). \nonumber
% \end{align}
% In LfO, the expert action is absent, thus directly working on the above is infeasible. Thus, one can alternatively choose to minimize the discrepancy on the state-only transition occupancy measure $\rho_{\pi_\theta}(\bs, \bs')$: 
% \begin{align}%\label{eq:min-il-lfo}
% 	D(\pi_\theta(\btau), \pi_E(\btau)) %&\to D(\textcolor{red}{f} \circ \pi_\theta(\btau), \pi_E(\tbtau)) \nonumber\\
% 	&\to D(\rho_{\pi_\theta}(\bs, \bs'), \rho_{\pi_E}(\bs, \bs')). \nonumber
% \end{align}
% %where $\textcolor{red}{f}$ is used to extract states/observations and drop out actions from the generated trajectories $\btau\sim\pi_\theta(\btau)$. 

Besides, we can also divide IL settings into two orthogonal categories: {{online IL}} and {{offline IL}}. 
In {{{online IL}}}, the learning policy $\pi_\theta$ can interact with the environment and generate online trajectories $\btau\sim\pi_\theta(\btau)$. % as specified in Equation~\ref{eq:generate-tau}. 
In {{{offline IL}}}, the agent cannot interact with the environment but has access to an offline static dataset $\{\btau |\btau \sim \pi_\beta(\btau)\}$, collected by some unknown (sub-optimal) behavior policies $\pi_\beta$. By leveraging the offline data $\{\pi_\beta(\btau)\} \cup \{\pi_E(\btau)\}$ without any interactions with the environment, the goal of offline IL is to learn a policy recovering the expert behaviors (demonstrations or observations) generated by $\pi_E$. Note that, in contrast to the typical offline RL problem~\citep{levine2020offline}, the offline data $\{\pi_\beta(\btau)\}$ in offline IL does not {contains} any reward signal. 

\textbf{Cross-domain IL.} Beyond the above two IL branches (online/offline and LfD/LfO), we can also divide IL into: 1)~single-domain IL and 2)~cross-domain IL, where 1)~the single-domain IL assumes that the expert behaviors are collected in the same MDP in which the learning policy is to be learned, and 2)~the cross-domain IL studies how to imitate expert behaviors when
discrepancies exist between the expert and the learning MDPs (\eg differing in their transition dynamics or morphologies). 





\subsection{Hindsight Information Matching} 
\label{sec:hindsight-infomation-matching}

In typical goal-conditioned RL problems, hindsight experience replay (HER)~\citep{andrychowicz2017hindsight} proposes to leverage the rich repository of the failed experiences by replacing the desired (true) goals of training trajectories with the achieved goals of the failed experiences:  
\begin{align}
	\texttt{Alg} (\pi_\theta; \bg, \btau_\bg) \to \texttt{Alg} (\pi_\theta; {f}_{\text{HER}}(\btau_\bg), \btau_\bg), \nonumber
\end{align}
where the learner $\texttt{Alg}(\pi_\theta; \cdot, \cdot)$ could be any RL methods, $\btau_\bg\sim\pi_\theta(\btau_\bg|\bg)$ denotes the trajectory generated by a goal-conditioned policy $\pi_\theta(\ba_t|\bs_t,\bg)$, and ${f}_{\text{HER}}$ denotes a pre-defined (hindsight information extraction) function, \eg returning the last state in trajectory $\btau_\bg$. 

HER can also be applied to the (single-goal) reward-driven online/offline RL tasks, setting the return (sum of the discounted rewards) of a trajectory as an implicit goal for the corresponding trajectory. 
Thus, \textit{we can reformulate the (single-goal) reward-driven RL task}, learning policy $\pi_\theta(\ba_t|\bs_t)$ that maximize the return, \textit{as a multi-goal RL task}, learning a return-conditioned policy $\pi_\theta(\ba_t|\bs_t, \cdot)$ that maximize the following log-likelihood: 
\begin{align}\label{eq:reward-conditioned-policy}
	\max_{\pi_\theta} \ \mathbb{E}_{\mathcal{D}(\btau)} \lbr \log \pi_\theta(\ba|\bs, {f_\text{R}}(\btau)) \rbr,
\end{align}
where ${f_\text{R}}(\btau)$ denotes the return of trajectory $\btau$. At test, we can then condition the contextual policy $\pi_\theta(\ba|\bs, \cdot)$ on a desired target return. 
In offline RL, the empirical distribution $\mathcal{D}(\btau)$ in Equation~\ref{eq:reward-conditioned-policy} can be naturally set as the offline data distribution; in online RL, $\mathcal{D}(\btau)$ can be set as the replay/experience buffer, and will be updated and biased towards trajectories that have high expected returns. 


Intuitively, biasing the sampling distribution ($\mathcal{D}(\btau)$ towards higher returns) leads to \textit{an implicit policy improvement operation}. However, such an operator is non-trivial to obtain in the IL problem, where we do not have access to a pre-defined function ${f_\text{R}}(\btau)$ to bias the learning policy towards recovering the given expert data $\{\pi_E(\btau)\}$ (demonstrations or observations). 

